\documentclass[11pt,a4paper]{article}

\usepackage[utf8]{inputenc}
\usepackage[T2A]{fontenc}
\usepackage[russian,english]{babel}
\usepackage{amsmath,amssymb,amsthm,mathtools}
\usepackage{geometry}
\usepackage{booktabs}
\usepackage{enumitem}
\usepackage{hyperref}

\geometry{margin=2.5cm}

\title{Toy-модель оптимального управления банком с процентным риском и капиталом}
\author{}
\date{}

\newcommand{\1}{\mathbf{1}}
\newcommand{\E}{\mathbb{E}}
\newcommand{\R}{\mathbb{R}}
\newcommand{\pos}[1]{\left[#1\right]^+}

\begin{document}
\selectlanguage{russian}
\maketitle

\section{Идея модели}

Рассматривается минимальная toy-модель банка, в которой есть только процентный риск и капитал. Ликвидностный риск, кредитный риск, досрочные погашения, рыночная переоценка, налоги и операционные расходы отсутствуют.

Банк имеет fixed-rate кредитный портфель и фондирует его floating-rate депозитами. Кредиты и депозиты считаются равными по объёму:
\begin{equation}
    A_t = D_t.
\end{equation}

Капитал интерпретируется как cash/reserve buffer:
\begin{equation}
    C_t = K_t.
\end{equation}

Баланс имеет вид:
\begin{equation}
    A_t + C_t = D_t + K_t.
\end{equation}

При $A_t=D_t$ и $C_t=K_t$ баланс выполняется автоматически. Экономически это означает, что кредиты полностью фондируются депозитами, а капитал лежит отдельно и покрывает неблагоприятные реализации процентного риска.

\section{Состояние, управление и шоки}

Состояние в момент $t$:
\begin{equation}
    S_t = (A_t, K_t, \overline{R}_t, r_t),
\end{equation}
где:
\begin{itemize}[noitemsep]
    \item $A_t$ --- объём кредитного портфеля;
    \item $K_t$ --- капитал / cash buffer;
    \item $\overline{R}_t$ --- средняя фиксированная ставка старого кредитного портфеля;
    \item $r_t$ --- short rate.
\end{itemize}

Управление:
\begin{equation}
    a_t = (g_t, p_t),
\end{equation}
где:
\begin{itemize}[noitemsep]
    \item $g_t$ --- скорость новой выдачи кредитов как доля текущего портфеля;
    \item $p_t \in [0,1]$ --- желаемый payout ratio от положительной прибыли.
\end{itemize}

Ставка $r_t$ развивается стохастически. Например, можно использовать AR(1)-процесс:
\begin{equation}
    r_{t+1} = \mu + \rho(r_t-\mu) + \sigma \varepsilon_{t+1},
    \qquad \varepsilon_{t+1}\sim\mathcal{N}(0,1).
\end{equation}

При необходимости можно добавить floor:
\begin{equation}
    r_{t+1} = \max\left(r_{\min},\, \mu + \rho(r_t-\mu)+\sigma\varepsilon_{t+1}\right).
\end{equation}

\section{Кредитный портфель}

Пусть $\alpha$ --- месячная амортизация старого портфеля. Тогда новая динамика баланса:
\begin{equation}
    A_{t+1} = A_t(1-\alpha+g_t).
\end{equation}

Если $g_t=\alpha$, то банк поддерживает constant balance. Если $g_t>\alpha$, баланс растёт. Если $g_t<\alpha$, баланс сокращается.

\subsection{Ставка новой выдачи}

Депозиты стоят ровно short rate:
\begin{equation}
    R_t^D = r_t.
\end{equation}

Новые кредиты выдаются под ставку:
\begin{equation}
    R_t = r_t + s(g_t),
\end{equation}
где спред линейно зависит от скорости выдачи:
\begin{equation}
    s(g_t) = s_0 - \kappa(g_t-\alpha), \qquad \kappa>0.
\end{equation}

При normal rolling, $g_t=\alpha$, спред равен нормальному коммерческому уровню:
\begin{equation}
    s(\alpha)=s_0.
\end{equation}

Если банк растёт быстрее rolling balance, $g_t>\alpha$, то спред ниже $s_0$: банк вынужден делать менее маржинальные сделки. Если банк растёт медленнее, $g_t<\alpha$, то спред выше $s_0$: банк выбирает только более маржинальные сделки.

При желании можно добавить floor:
\begin{equation}
    s(g_t)=\max\{s_{\min},\,s_0-\kappa(g_t-\alpha)\}.
\end{equation}

\subsection{Средняя ставка портфеля}

Средняя фиксированная ставка кредитного портфеля обновляется как взвешенное среднее старого surviving book и новой выдачи:
\begin{equation}
    \overline{R}_{t+1}
    =
    \frac{(1-\alpha)\overline{R}_t + g_t\left[r_t+s_0-\kappa(g_t-\alpha)\right]}
    {1-\alpha+g_t}.
\end{equation}

Эквивалентно:
\begin{equation}
    \overline{R}_{t+1}
    =
    \frac{(1-\alpha)A_t\overline{R}_t + g_tA_t R_t}
    {A_t(1-\alpha+g_t)}.
\end{equation}

\section{Процентный доход и капитал}

Так как $D_t=A_t$, стоимость депозитов равна $r_tA_t$. Капитал лежит в cash и зарабатывает $r_tK_t$. Чистый процентный доход:
\begin{equation}
    \Pi_t = NII_t = A_t(\overline{R}_t-r_t)+r_tK_t.
\end{equation}

Преддивидендный капитал:
\begin{equation}
    K_t^{pre}=K_t+\Pi_t.
\end{equation}

В данной toy-модели весь P\&L проходит через капитал. Если $\Pi_t<0$, то капитал уменьшается. Это и есть механизм покрытия процентных потерь.

\section{Капитальные пороги, дивиденды и дефолт}

Capital ratio (capital adequacy ratio) обозначается как:
\begin{equation}
    c_t = CAR_t = \frac{K_t}{A_t}.
\end{equation}

Вводятся два уровня капитала:
\begin{equation}
    c_0 < c^*,
\end{equation}
где:
\begin{itemize}[noitemsep]
    \item $c^*$ --- dividend / management capital threshold;
    \item $c_0$ --- regulatory default / resolution threshold.
\end{itemize}

\subsection{Default rule}

Проверка дефолта делается после прибыли текущего периода и с учётом нового размера баланса:
\begin{equation}
    \frac{K_t^{pre}}{A_{t+1}} < c_0
    \quad\Rightarrow\quad
    \text{default/resolution at time } t.
\end{equation}

Если происходит дефолт, игра заканчивается. Акционер получает discounted residual capital:
\begin{equation}
    Rec_t = (1-\delta)\pos{K_t^{pre}},
\end{equation}
где $\delta\in[0,1]$ --- haircut/dislocation discount в resolution.

\subsection{Dividend rule}

Если дефолта нет, дивиденды выплачиваются только из положительной прибыли и только сверх целевого capital threshold. Дивиденд не может опустить банк ниже $c^*$ относительно нового баланса $A_{t+1}$:
\begin{equation}
    Div_t
    =
    \min\left(
        p_t\pos{\Pi_t},
        \pos{K_t^{pre}-c^*A_{t+1}}
    \right).
\end{equation}

Если $K_t^{pre}/A_{t+1}<c^*$, то второй аргумент равен нулю, и дивиденды запрещены:
\begin{equation}
    Div_t=0.
\end{equation}

Капитал после выплаты дивидендов:
\begin{equation}
    K_{t+1}=K_t^{pre}-Div_t.
\end{equation}

\section{Payoff акционера}

Пусть $\tau$ --- момент дефолта:
\begin{equation}
    \tau = \inf\left\{t: \frac{K_t^{pre}}{A_{t+1}}<c_0\right\}.
\end{equation}

Если банк доживает до горизонта $T$, terminal value задаётся как:
\begin{equation}
    TV_T = K_T + \theta A_T,
\end{equation}
где $\theta A_T$ --- упрощённая franchise value кредитного портфеля.

Целевая функция:
\begin{equation}
\boxed{
    \max_{\pi}
    \E\left[
        \sum_{t<\tau,\,t\le T} \gamma^t Div_t
        +
        \1_{\tau\le T}\gamma^\tau (1-\delta)\pos{K_\tau^{pre}}
        +
        \1_{\tau>T}\gamma^T(K_T+\theta A_T)
    \right]
}
\end{equation}
где политика
\begin{equation}
    \pi: (A_t,K_t,\overline{R}_t,r_t)\mapsto (g_t,p_t).
\end{equation}

\section{Полная рекурсия модели}

Для каждого периода $t$:

\begin{enumerate}[label=\arabic*.]
    \item Наблюдается состояние $S_t=(A_t,K_t,\overline{R}_t,r_t)$.
    \item Выбирается управление $a_t=(g_t,p_t)$.
    \item Считается новый размер баланса:
    \begin{equation}
        A_{t+1}=A_t(1-\alpha+g_t).
    \end{equation}
    \item Считается прибыль текущего периода:
    \begin{equation}
        \Pi_t=A_t(\overline{R}_t-r_t)+r_tK_t.
    \end{equation}
    \item Считается преддивидендный капитал:
    \begin{equation}
        K_t^{pre}=K_t+\Pi_t.
    \end{equation}
    \item Если
    \begin{equation}
        K_t^{pre}/A_{t+1}<c_0,
    \end{equation}
    то происходит дефолт, акционер получает
    \begin{equation}
        (1-\delta)\pos{K_t^{pre}},
    \end{equation}
    и процесс останавливается.
    \item Иначе дивиденды равны:
    \begin{equation}
        Div_t=\min\left(p_t\pos{\Pi_t},\pos{K_t^{pre}-c^*A_{t+1}}\right).
    \end{equation}
    \item Капитал обновляется:
    \begin{equation}
        K_{t+1}=K_t^{pre}-Div_t.
    \end{equation}
    \item Средний купон обновляется:
    \begin{equation}
        \overline{R}_{t+1}
        =
        \frac{(1-\alpha)\overline{R}_t + g_t\left[r_t+s_0-\kappa(g_t-\alpha)\right]}
        {1-\alpha+g_t}.
    \end{equation}
    \item Ставка обновляется:
    \begin{equation}
        r_{t+1}=\mu+\rho(r_t-\mu)+\sigma\varepsilon_{t+1}.
    \end{equation}
\end{enumerate}



\section{Редукция размерности через однородность}

В исходной постановке состояние имеет размерность четыре:
\begin{equation}
    S_t=(A_t,K_t,\overline{R}_t,r_t).
\end{equation}

Однако модель обладает однородностью по масштабу баланса. Все денежные величины --- прибыль, дивиденды, капитал, recovery и terminal value --- линейно масштабируются с $A_t$. Поэтому можно ввести относительный капитал:
\begin{equation}
    c_t=\frac{K_t}{A_t},
\end{equation}
и переписать задачу в нормированных переменных. Тогда состояние становится трёхмерным:
\begin{equation}
    \widetilde S_t=(c_t,\overline{R}_t,r_t).
\end{equation}

\subsection{Нормированные переменные}

Обозначим множитель роста баланса:
\begin{equation}
    m_t=1-\alpha+g_t.
\end{equation}
Тогда:
\begin{equation}
    A_{t+1}=A_tm_t.
\end{equation}

Прибыль текущего периода равна:
\begin{equation}
    \Pi_t=A_t(\overline{R}_t-r_t)+r_tK_t.
\end{equation}
Подставляя $K_t=c_tA_t$, получаем:
\begin{equation}
    \Pi_t=A_t\left(\overline{R}_t-r_t+r_tc_t\right).
\end{equation}
Введём прибыль на единицу активов:
\begin{equation}
    \pi_t=\frac{\Pi_t}{A_t}=\overline{R}_t-r_t+r_tc_t.
\end{equation}

Преддивидендный капитал в относительных величинах:
\begin{equation}
    \frac{K_t^{pre}}{A_t}=c_t+\pi_t.
\end{equation}

\subsection{Default rule в нормированной форме}

В полной модели дефолт происходит, если:
\begin{equation}
    \frac{K_t^{pre}}{A_{t+1}}<c_0.
\end{equation}
Так как $A_{t+1}=A_tm_t$ и $K_t^{pre}=A_t(c_t+\pi_t)$, условие дефолта принимает вид:
\begin{equation}
    \boxed{
    \frac{c_t+\pi_t}{m_t}<c_0.
    }
\end{equation}

Если дефолт происходит, акционер получает recovery:
\begin{equation}
    Rec_t=(1-\delta)\pos{K_t^{pre}}.
\end{equation}
Нормированный recovery на единицу активов:
\begin{equation}
    \boxed{
    rec_t=(1-\delta)\pos{c_t+\pi_t}.
    }
\end{equation}

\subsection{Dividend rule в нормированной форме}

Полный дивиденд:
\begin{equation}
    Div_t=\min\left(p_t\pos{\Pi_t},\pos{K_t^{pre}-c^*A_{t+1}}\right).
\end{equation}
Делим на $A_t$:
\begin{equation}
    \frac{Div_t}{A_t}
    =
    \min\left(
        p_t\pos{\pi_t},
        \pos{c_t+\pi_t-c^*m_t}
    \right).
\end{equation}
Введём нормированный дивиденд:
\begin{equation}
    \boxed{
    d_t
    =
    \min\left(
        p_t\pos{\pi_t},
        \pos{c_t+\pi_t-c^*m_t}
    \right).
    }
\end{equation}

Если $c_t+\pi_t<c^*m_t$, то $d_t=0$. Это ровно означает, что после прибыли текущего периода банк не имеет dividend capacity относительно нового размера баланса.

\subsection{Динамика относительного капитала}

Если дефолта нет, капитал после дивидендов:
\begin{equation}
    K_{t+1}=K_t+\Pi_t-Div_t.
\end{equation}
Делим на $A_{t+1}=A_tm_t$:
\begin{equation}
    c_{t+1}
    =
    \frac{K_{t+1}}{A_{t+1}}
    =
    \frac{A_tc_t+A_t\pi_t-A_td_t}{A_tm_t}.
\end{equation}
Следовательно:
\begin{equation}
    \boxed{
    c_{t+1}=\frac{c_t+\pi_t-d_t}{m_t}.
    }
\end{equation}

\subsection{Динамика средней ставки портфеля и short rate}

Средняя ставка кредитного портфеля не зависит от масштаба $A_t$:
\begin{equation}
    \boxed{
    \overline{R}_{t+1}
    =
    \frac{(1-\alpha)\overline{R}_t+g_t\left[r_t+s_0-\kappa(g_t-\alpha)\right]}{m_t}.
    }
\end{equation}

Ставочный процесс также не зависит от $A_t$:
\begin{equation}
    \boxed{
    r_{t+1}=\mu+\rho(r_t-\mu)+\sigma\varepsilon_{t+1}.
    }
\end{equation}

\subsection{Нормированная рекурсия}

Для фиксированного управления $(g_t,p_t)$ нормированная динамика имеет вид:
\begin{align}
    m_t &= 1-\alpha+g_t,\\
    \pi_t &= \overline{R}_t-r_t+r_tc_t,\\
    d_t &= \min\left(p_t\pos{\pi_t},\pos{c_t+\pi_t-c^*m_t}\right),\\
    c_{t+1} &= \frac{c_t+\pi_t-d_t}{m_t},\\
    \overline{R}_{t+1} &= \frac{(1-\alpha)\overline{R}_t+g_t\left[r_t+s_0-\kappa(g_t-\alpha)\right]}{m_t},\\
    r_{t+1} &= \mu+\rho(r_t-\mu)+\sigma\varepsilon_{t+1}.
\end{align}

Default indicator:
\begin{equation}
    \mathcal{D}_t=
    \1\left\{\frac{c_t+\pi_t}{m_t}<c_0\right\}.
\end{equation}

Таким образом, вместо состояния $(A_t,K_t,\overline{R}_t,r_t)$ можно использовать $(c_t,\overline{R}_t,r_t)$.

\section{Нормированное уравнение Беллмана}

Пусть $V_t(A,K,\overline{R},r)$ --- value function акционера в абсолютных денежных единицах. Из-за однородности модели по $A$ можно искать её в виде:
\begin{equation}
    \boxed{
    V_t(A,K,\overline{R},r)=A\,v_t(c,\overline{R},r),
    \qquad c=\frac{K}{A}.
    }
\end{equation}
Здесь $v_t(c,\overline{R},r)$ --- стоимость банка на единицу кредитного портфеля.

\subsection{Terminal condition}

Если банк дожил до горизонта $T$, terminal value равен:
\begin{equation}
    TV_T=K_T+\theta A_T.
\end{equation}
Поэтому нормированное терминальное условие:
\begin{equation}
    \boxed{
    v_T(c,\overline{R},r)=c+\theta.
    }
\end{equation}

\subsection{Bellman equation}

Для состояния $(c,\overline{R},r)$ и управления $(g,p)$ определим:
\begin{align}
    m &= 1-\alpha+g,\\
    \pi &= \overline{R}-r+rc,\\
    R &= r+s_0-\kappa(g-\alpha),\\
    \overline{R}' &= \frac{(1-\alpha)\overline{R}+gR}{m},\\
    d &= \min\left(p\pos{\pi},\pos{c+\pi-c^*m}\right),\\
    c' &= \frac{c+\pi-d}{m},\\
    r' &= \mu+\rho(r-\mu)+\sigma\varepsilon'.
\end{align}

Если:
\begin{equation}
    \frac{c+\pi}{m}<c_0,
\end{equation}
то наступает default/resolution, и нормированный payoff равен:
\begin{equation}
    (1-\delta)\pos{c+\pi}.
\end{equation}

Если дефолта нет, Bellman continuation value в абсолютных единицах:
\begin{equation}
    Div_t+\gamma\E_t[V_{t+1}(A_{t+1},K_{t+1},\overline{R}_{t+1},r_{t+1})].
\end{equation}
Так как $Div_t=Ad$, $A_{t+1}=Am$ и $V_{t+1}=A_{t+1}v_{t+1}$, после деления на $A$ получаем множитель $m$ перед continuation value:
\begin{equation}
    d+\gamma m\E\left[v_{t+1}(c',\overline{R}',r')\right].
\end{equation}

Итоговое нормированное уравнение Беллмана:
\begin{equation}
\boxed{
    v_t(c,\overline{R},r)
    =
    \max_{g\in\mathcal{G},\,p\in\mathcal{P}}
    \begin{cases}
        (1-\delta)\pos{c+\pi},
        & \dfrac{c+\pi}{m}<c_0,\\[8pt]
        d+\gamma m\E\left[v_{t+1}(c',\overline{R}',r')\right],
        & \dfrac{c+\pi}{m}\ge c_0.
    \end{cases}
}
\end{equation}

Эквивалентно, через индикаторы:
\begin{equation}
\boxed{
\begin{aligned}
    v_t(c,\overline{R},r)
    =
    \max_{g,p}\Bigg\{&
    \1\left\{\frac{c+\pi}{m}<c_0\right\}(1-\delta)\pos{c+\pi}
    \\
    &+
    \1\left\{\frac{c+\pi}{m}\ge c_0\right\}
    \left(d+\gamma m\E\left[v_{t+1}(c',\overline{R}',r')\right]\right)
    \Bigg\}.
\end{aligned}
}
\end{equation}

\subsection{Экономический смысл редукции}

Редукция показывает, что абсолютный размер кредитного портфеля $A_t$ сам по себе не нужен для выбора оптимальной политики, если все доходы, потери и выплаты масштабируются линейно с $A_t$. Важны только:
\begin{itemize}[noitemsep]
    \item относительный капитал $c_t=K_t/A_t$;
    \item средняя ставка старого fixed-rate портфеля $\overline{R}_t$;
    \item текущая short rate $r_t$.
\end{itemize}

Рост $g_t$ влияет на задачу тремя каналами:
\begin{enumerate}[noitemsep]
    \item увеличивает будущий масштаб через множитель $m_t=1-\alpha+g_t$;
    \item размывает капитал, так как $c_{t+1}=(c_t+\pi_t-d_t)/m_t$;
    \item снижает спред новой выдачи через $s(g_t)=s_0-\kappa(g_t-\alpha)$.
\end{enumerate}

Поэтому оптимизатор выбирает компромисс между ростом franchise value, устойчивостью capital ratio и ухудшением маржинальности новой выдачи.

\subsection{Когда однородность ломается}

Редукция $S_t=(A_t,K_t,\overline{R}_t,r_t)\mapsto(c_t,\overline{R}_t,r_t)$ корректна, пока модель линейно масштабируется с $A_t$. Она перестанет быть точной, если добавить:
\begin{itemize}[noitemsep]
    \item фиксированные операционные расходы;
    \item ограничения абсолютного размера баланса $A_t\le A_{\max}$;
    \item нелинейную зависимость спреда от абсолютного объёма $A_t$;
    \item внешнюю докапитализацию с фиксированными costs;
    \item налоги с нелинейной шкалой;
    \item минимальные дивиденды в абсолютных денежных единицах;
    \item нелинейные regulatory floors, зависящие не только от отношения $K_t/A_t$.
\end{itemize}

В текущей toy-постановке таких элементов нет, поэтому нормированная трёхмерная задача эквивалентна исходной четырёхмерной задаче.

\section{Экономическая интерпретация}

Модель содержит базовый trade-off:
\begin{itemize}
    \item быстрый рост $g_t$ увеличивает будущий портфель $A_{t+1}$;
    \item быстрый рост снижает спред новой выдачи $s(g_t)$;
    \item рост баланса увеличивает знаменатель $CAR=K/A$;
    \item fixed-rate активы имеют медленную среднюю ставку $\overline{R}_t$, а депозитная стоимость равна текущему short rate $r_t$;
    \item рост ставок снижает $NII_t=A_t(\overline{R}_t-r_t)+r_tK_t$ и может проедать капитал;
    \item если capital ratio ниже $c^*$, дивиденды прекращаются;
    \item если capital ratio ниже $c_0$, банк попадает в default/resolution, и акционер получает recovery с дисконтом.
\end{itemize}

\section{Возможные расширения}

Натуральные расширения модели:
\begin{enumerate}[noitemsep]
    \item добавить ограничение или cost на отрицательный/слишком малый $g_t$;
    \item добавить convexity/floor в функцию $s(g_t)$;
    \item заменить одну среднюю ставку $\overline{R}_t$ на maturity ladder или 2D-gap;
    \item добавить налоги и операционные расходы;
    \item добавить credit losses;
    \item добавить ликвидность и оттоки депозитов;
    \item добавить докапитализацию вместо immediate default;
    \item добавить stochastic deposit beta вместо $R_t^D=r_t$;
    \item добавить risk-sensitive objective, например CVaR terminal shareholder loss.
\end{enumerate}

\end{document}
